% !TeX root = ../papers/nc-p3.tex
\section{混合点を横断する曲線の全同型と周期}
\label{sec:marked-mixed-curve-isomorphisms}
\begin{theorem}\label{thm:marked-mixed-curve-isomorphisms}
命題\ref{prop:split-mixed-transverse-curve}の標準正則な二窓 $W(t),W(t')$ が
印を忘れて添字保持同型である必要十分条件は $t'^2=t^2$ である。
最小周期は $t^2=1,-1$ で1、他で2である。特に $x=t^2$ を同型類のパラメータとすれば、
除かれる値は $x=-3,-1/3$ の二つである。階数14の $x=-1$ は正則なWalsh窓とも同型にならない。
$x\notin\{0,1,-1\}$ の正則点で四位置に作用する射影自己同型群の位数は48である。
\end{theorem}

\begin{proof}[反対切断の逆から軸を復元する]
反対階数16の四テンソル $w$ に対して、行 $ac$・列 $bd$ の反対切断の逆転置を
同じ添字順に戻した共変テンソルを $w^\vee$ とする。異なる二位置 $i,j$ を残し、
$w\otimes w^\vee$ の他二位置を縮約して
$C_{ij}\in\mathrm{End}(V_i)\otimes\mathrm{End}(V_j)$ を取る。
例えば $C_{01;aa',bb'}=\sum_{c,d}w_{abcd}w^\vee_{a'b'cd}$ である。
両端の跡対合によりこれを写像 $L_{ij}:\mathrm{End}(V_j)\to\mathrm{End}(V_i)$ とみて
$K_{ij}=L_{ij}L_{ij}^*$ を作る。全体倍率は消え、基底変換は各位置で共役になる。
反対切断を転置しても逆テンソルは添字交換で対応するので、この構成は巡回移動とも可換である。

位置 $i$ の二隣 $j=i\pm1$ について $K_{ij}$ の固有値 $16/9$ の同時固有空間を取る。
$t(t^2+1)(t^2-1)\ne0$ なら、$i=0,1$ の少なくとも一方でその空間は
跡零対角空間に等しく、したがってその四共通固有直線が四座標軸を復元する。
これを全ての例外根も含めて示す。各行列単位 $E_{a,a+h}$ を同色 $h$ でまとめると
$C_{ij}$ は四ブロックで、$R_h$ をその $h$ ブロック、$S_h$ を $a\mapsto a+h$ の置換とすると
$K_{ij}$ のブロックは $R_h S_h R_h^t S_h$ である。
$K_{ij}-16I/9$ に $144t^2(1+t^2)^2$ を掛けると、対角ブロックは
全成分が $512t^2(1+t^2)^2$ の行列である。色 $h=1$ では二隣を $+,-$ の順に縦に重ね、
零始まりの行番号 $(0,1,2,3)$ と $(0,1,2,4)$ を取る。
$i=0,1$ の計四小行列式のモニック最大公約数は $(t^2-1)^4(t^2+1)^4$ である。
色 $1,2,3$ の巡回は元の三列を保つため、単元になる小行列のある位置では
三つの非零色の全てが核に寄与しない。この検算は
\path{computations/verify_mixed_curve_axes.py}の有理関数体上の演算である。
$t^2=1$ でも跡零対角空間の包含は保たれる。
従ってこの方法で四軸を復元できる窓との同型があれば、相手の対応位置でも
同じ次元3の空間を保ち、やはり四軸を移す。

$u=1+t$, $v=t-1$ とする。各巡回三組 $(i,j,k)$ で、辺の色 $i$ の列を
$(0i,i0,jk,kj)$ とした二行の関係は
\[
 E(u,v)=\begin{pmatrix}1&0&u&-v\\0&1&v&u\end{pmatrix}
\]
である。位置 $0,1,2,3$ から順にその引数は $(u,v),(v,-u),(-u,-v),(-v,u)$ となる。
どの軸も点グラフで両隣の同名の軸だけに対応する。実際 $0$ を評価すると補三文字が消え、
$i\ne0$ を評価すると $0$ と他の二文字が消える。これは $u,v$ が同時に零でないことによる。
従って一位置から四位置全てに軸を運べて、同型の置換部分は全位置で同じである。
\end{proof}

\begin{proof}[対角同値と特殊値]
Walsh表示の同時平行移動は四軸上で推移的であり、前節の規格化は対角である。
$t^2\ne-1$ の階数16なら、相手側でこの自己同型を合成して $0$ 軸を固定できる。
残る三軸の巡回置換は窓を保つ。奇置換を全位置で施し、
奇位置の補文字をさらに $-1$ 倍すると $W(t)$ は $W(-t)$ になる。
これは $D$ が置換で符号を変えず、外積 $K$ は行列式の符号を取るためで、
階数の制限なしに成り立つ。

よって $s=t$ または $-t$ とし、$E(u_s,v_s)$ を対角倍率だけで移せる場合を調べればよい。
星倍率を各位置で1に戻し、二位置の補文字の倍率を $a_i,b_i$ とする。
行の先頭二列が単位行列なので、$E(u',v')$ に移すには、全ての巡回三組で
\[
 a_j b_k=(u'/u_s)b_i=(v'/v_s)a_i,\qquad
 a_k b_j=(v'/v_s)b_i=(u'/u_s)a_i
\]
を要する。ここでは四係数が非零の場合を扱う。
$a_i/b_i$ が全 $i$ で一定なので左の二つの積は等しく、
$u'/u_s=v'/v_s$ となる。相手が $W(t')$ の位置0ならこれで $t'=s$ である。
相手が $W(t)$ の位置1なら $(u',v')=(v_t,-u_t)$ であり、$s=t$ では
$v_t^2=-u_t^2$ が必要、$s=-t$ では $-1=1$ を要する。
従って対象の階数16で両係数が非零なら最小周期は2である。
いずれかの係数だけが零になる値と両方非零の値の同一視は、一位置の軸の復元後、
上の関係行列の零成分の不一致で排除される。

自己同型ではこの範囲で $s=t$ なので、星を固定する軸置換は巡回三通りである。
置換が恒等なら上式で $a_i=b_i$、$a_ja_k=a_i$ であり、解は積1の三つの符号の四通りである。
隣接位置にも同じ式を適用して全位置の倍率が揃う。これら四変換は実際に全関係を保ち、
星の像の四通りも前記の平行移動で実現する。従って群位数は $4\cdot3\cdot4=48$ である。

$t=1$ では、三補文字の一つの互換を全位置に施し、奇位置の星文字のみ符号を変えると、
上の四つの関係空間は各々一つ先に移る。$t=-1$ はこれと同型である。
$t=0$ は固有の階数10で最小周期2、$t^2=-1$ は固有の階数14で二値は互いに同型である。
後者は一周期の表示も明示できる。位置 $n$ の星文字を $(1+t)t^n z$ に置き、
三補文字を $x_1,x_2,x_3$ と共通に読むと、各位置で
\[
          x_i z=z x_i,\qquad x_jx_k-tx_kx_j=-z x_i
              \quad ((i,j,k)\ \text{巡回})
\]
になる。実際 $v=tu$ で、各辺の第一行の二つの末尾係数は $(ut^n,-ut^{n+1})$ であり、
第二行と第一行の $t$ 倍の差が可換関係を与える。

最後に $t^2=-1$ の四テンソルに、全位置で星を $z^3$ 倍、三補文字を $z^{-1}$ 倍する
特殊線型群の一径数部分群を作用させる。前節の一つ星の極限を得る。その反対切断では
零色のブロックは零、他は四次可逆なので階数12である。従って元のアフィン軌道は閉でない。
Walsh表示の非零テンソルの軌道はノルム最小性により閉である
（定理\ref{thm:walsh-spectral-fibers}の証明）。倍率と任意の基底変更でこの相違は保たれる。
\end{proof}
